% =====================================================================
% Three-Way Comparison Report
% LongRAG vs Metadata-RAG vs My_RAG on Four Heterogeneous Corpora
% =====================================================================
\documentclass[11pt,a4paper]{article}

\usepackage[utf8]{inputenc}
\usepackage[T1]{fontenc}
\usepackage{lmodern}
\usepackage[margin=1in]{geometry}
\usepackage{amsmath,amssymb,amsfonts}
\usepackage{booktabs}
\usepackage{multirow}
\usepackage{array}
\usepackage{tabularx}
\usepackage{siunitx}
\usepackage{graphicx}
\usepackage{xcolor}
\usepackage{hyperref}
\usepackage{enumitem}
\usepackage{fancyhdr}
\usepackage{titlesec}
\usepackage{caption}
\usepackage{microtype}
\usepackage{parskip}

\hypersetup{
  colorlinks=true,
  linkcolor=blue!50!black,
  citecolor=blue!50!black,
  urlcolor=blue!50!black,
  pdftitle={Three-Way Comparison Report},
}

\sisetup{
  round-mode=places,
  round-precision=4,
  detect-weight=true,
  detect-family=true,
}

\titleformat{\section}{\large\bfseries}{\thesection}{1em}{}
\titleformat{\subsection}{\normalsize\bfseries}{\thesubsection}{1em}{}

\pagestyle{fancy}
\fancyhf{}
\rhead{\small Three-Way Comparison}
\lhead{\small LongRAG $\cdot$ Metadata-RAG $\cdot$ My\_RAG}
\cfoot{\thepage}

\title{
  \vspace{-1.5cm}
  \rule{\linewidth}{0.4pt}\\[0.3cm]
  {\Large\bfseries LongRAG \emph{vs} Metadata-RAG \emph{vs}
   My\_RAG:}\\[0.15cm]
  {\large A Retrieval-Only Three-Way Comparison on Four\\
   Heterogeneous Long-Context Corpora}\\[0.2cm]
  \rule{\linewidth}{0.4pt}\\[0.4cm]
  \normalsize RAG Architecture Readiness Program\\
  \small Comparative Report
}
\author{}
\date{\today}

\begin{document}

\maketitle
\thispagestyle{empty}

\begin{abstract}
\noindent
We compare three retrieval architectures under a fully controlled,
retrieval-only evaluation protocol on four heterogeneous
long-context corpora (CLERC, CUAD, GreekBarBench, LaRA). LongRAG is
a controlled adaptation of an external architecture that groups
adjacent documents into long retrieval units and retrieves with
dense similarity. Metadata-RAG is a paper-derived comparator whose
defining surface is a metadata-driven multistage filter preceding a
dense similarity search; under the frozen configuration its metadata
filter activates only on the LaRA corpus and it is therefore
measured as a dense-only retriever on the other three. My\_RAG is
the project-owned reference architecture, a structure-aware
hierarchical hybrid retriever that fuses dense and learned-sparse
signals by Reciprocal Rank Fusion and aggregates child chunks to
canonical documents through a parent-based rule. Under the shared
protocol My\_RAG achieves a macro-averaged document Recall@10 of
\num{0.4049}, Metadata-RAG of \num{0.3651}, and LongRAG of
\num{0.3144}. Metadata-RAG wins over LongRAG by \num{16.1}\% and
loses to My\_RAG by \num{9.8}\%. We report the full metric table,
the per-dataset decomposition, and the architectural observations
that emerge from the three-track ordering.
\end{abstract}

\vfill
\noindent\textbf{Keywords.} comparative retrieval; retrieval-only
benchmarking; metadata-driven retrieval; hybrid retrieval; dense
retrieval; controlled architectural comparison.

\newpage
\tableofcontents
\newpage

% =====================================================================
\section{Introduction}

\subsection{Question}
Under identical embedder, query set, ground truth, and evaluator,
which of three architecturally distinct retrievers returns the most
relevant canonical documents from long-context corpora? This report
answers the question for the three tracks that have been executed on
the program's four-dataset scope.

\subsection{Scope}
The four datasets are CLERC, CUAD, GreekBarBench, and LaRA. They
span the realistic range of source-document sizes from short
passages (CLERC, a few hundred tokens per document) to very long
documents (LaRA, hundreds of thousands of tokens per document).
The shared protocol, the shared embedder, and the shared evaluation
are described in Section~\ref{sec:protocol}.

\subsection{What this report does not claim}
It does not claim statistical significance. Each track is executed
once, deterministically, on its full query set. The observed
differences are descriptive. It does not claim native-fidelity
reproduction of any external architecture. It does not compare
end-to-end answer quality; all three tracks are retrieval-only.

% =====================================================================
\section{The Three Architectures}\label{sec:arch}

\subsection{LongRAG}
LongRAG reduces the retrieval corpus by grouping $K$ adjacent
documents into a single long retrieval unit whose token count is
bounded above by an architectural target of approximately
\num{4096} tokens. The paper's original setting uses Wikipedia
hyperlink adjacency to define ``adjacent''; this study uses a
deterministic adjacent-document grouping that preserves source
order. Retrieval is dense-only over the long-unit corpus. Retrieved
units are expanded to canonical documents by iterating through the
units in rank order and emitting each unseen document once.
LongRAG is executed as a controlled adaptation; the
architecture-defining components (long-unit grouping and dense-only
retrieval) are preserved.

\subsection{Metadata-RAG}
Metadata-RAG is a metadata-driven multistage retrieval
architecture. Its defining surface is a metadata layer that
intervenes between the query and the semantic similarity search.
Each canonical document (or each window of it, after amendment 04)
carries a typed metadata record extracted by a deterministic
extractor. Each query emits typed constraints produced by the same
extractor applied to the query text. Candidate windows must satisfy
a minimum match ratio of \num{0.5} against the query's emitted
constraints. The candidate set is then retrieved by exact dense
similarity over BGE-M3 embeddings, and the metadata and semantic
scores are fused with fixed coefficients $\alpha_m = 0.3$ and
$\alpha_s = 0.7$. When the candidate set is empty the pipeline
falls back to dense-only retrieval over the full corpus.

Under the frozen \texttt{min\_match\_ratio} of \num{0.5} the filter
activates only on LaRA (see Section~\ref{sec:analysis}). On CLERC,
CUAD, and GreekBarBench the architecture is measured as a dense-only
BGE-M3 retriever over the full corpus.

\subsection{My\_RAG}
My\_RAG is the project-owned reference architecture. It is a
structure-aware hierarchical hybrid retriever that commits to four
design principles: preserve source structure, retrieve fine and
rank coarse, combine complementary retrieval signals, and be
deterministic everywhere. It segments documents into
semantic-parent spans, chunks each parent into overlapping
\num{256}-token child windows without ever crossing parent
boundaries, and indexes each child chunk with both a dense
(\num{1024}-dimensional, L2-normalized BGE-M3) and a learned-sparse
(BGE-M3 lexical weight) representation. At query time it retrieves
independently from each representation, fuses the two rankings with
Reciprocal Rank Fusion ($k=60$), aggregates child chunks to their
parents with a best-plus-margin rule ($w=0.05$), and resolves
parents to canonical documents by best-parent-score.

\subsection{Architectural contrast}
Table~\ref{tab:arch} summarises the axes along which the three
architectures differ.

\begin{table}[h]
\centering
\caption{Architectural contrast across the three tracks.}
\label{tab:arch}
\small
\begin{tabular}{p{3.0cm} p{3.8cm} p{3.8cm} p{3.8cm}}
\toprule
\textbf{Axis} & \textbf{LongRAG} & \textbf{Metadata-RAG}
  & \textbf{My\_RAG} \\
\midrule
Unit type & long unit & single-doc window & child chunk \\
Unit size & $\approx\num{4096}$ tokens & $\le\num{8190}$ tokens
  & $\num{256}$ tokens \\
Lineage & multi-document & single document
  & single parent, single doc \\
Retrieval signal & dense & dense + metadata filter
  & dense + learned sparse \\
Fusion & none & weighted score & RRF ($k=60$) \\
Aggregation & first-occurrence by unit rank
  & first-occurrence by lineage id
  & best-plus-margin parent, best parent by doc \\
Corpus size (4 datasets) & \num{5669} units
  & \num{5725} windows & \num{101814} chunks \\
\bottomrule
\end{tabular}
\end{table}

% =====================================================================
\section{The Shared Protocol}\label{sec:protocol}

\subsection{What is shared}
Table~\ref{tab:shared} lists the components that are identical
across all three tracks.

\begin{table}[h]
\centering
\caption{Components shared identically by the three tracks.}
\label{tab:shared}
\small
\begin{tabular}{p{5.0cm} p{10.0cm}}
\toprule
\textbf{Component} & \textbf{Shared value} \\
\midrule
Datasets & CLERC, CUAD, GreekBarBench, LaRA \\
Canonical envelope & \texttt{documents.jsonl},
  \texttt{examples.jsonl}, \texttt{manifest.json} \\
Query set & identical \texttt{query\_id}s per dataset \\
Ground truth & canonical \texttt{ground\_truth.relevant\_ids} \\
Embedder identity & \texttt{BAAI/bge-m3} @ \texttt{5617a9f6\ldots} \\
Embedding precision & FP16 \\
Document embedding device & NVIDIA T4 (Google Colab) \\
Evaluator & document-level metrics, cutoffs $\{1,5,10,20,50\}$ \\
Primary metric & macro-averaged document \texttt{Recall@10} \\
Retrieval-only boundary & no generation, no LLM judge \\
\bottomrule
\end{tabular}
\end{table}

\subsection{What legitimately differs}
Unit type and unit size; retrieval signal; fusion mechanism;
aggregation rule. These are exactly the architectural axes that the
comparison is designed to measure.

\subsection{Retrieval-only boundary}
Every stage of every pipeline is strictly retrieval-only. Ground
truth is loaded only inside the evaluation stage, which is an
independent process from the retriever. The retriever receives only
\texttt{(query\_id, query\_text)} pairs.

% =====================================================================
\section{Results}\label{sec:results}

\subsection{Primary metric}
Table~\ref{tab:primary} reports the per-dataset and macro-averaged
document \texttt{Recall@10}.

\begin{table}[h]
\centering
\caption{Per-dataset and macro-averaged document \texttt{Recall@10}.}
\label{tab:primary}
\small
\begin{tabular}{l r r r}
\toprule
\textbf{Dataset} & \textbf{LongRAG} & \textbf{Metadata-RAG}
  & \textbf{My\_RAG} \\
\midrule
CLERC              & \num{0.0786} & \num{0.2726} & \num{0.3546} \\
CUAD               & \num{0.0324} & \num{0.0294} & \num{0.0431} \\
GreekBarBench      & \num{0.2606} & \num{0.2711} & \num{0.2852} \\
LaRA               & \num{0.8861} & \num{0.8874} & \num{0.9368} \\
\midrule
\textbf{Macro}     & \num{0.3144} & \num{0.3651} & \num{0.4049} \\
\bottomrule
\end{tabular}
\end{table}

\subsection{Full macro metric comparison}
Table~\ref{tab:full} reports the complete macro-averaged metric set.

\begin{table}[h]
\centering
\caption{Macro metric comparison across the three tracks.}
\label{tab:full}
\small
\begin{tabular}{l r r r}
\toprule
\textbf{Metric} & \textbf{LongRAG} & \textbf{Metadata-RAG}
  & \textbf{My\_RAG} \\
\midrule
Hit@1     & \num{0.1957} & \num{0.2313} & \num{0.2553} \\
Hit@5     & \num{0.2606} & \num{0.3389} & \num{0.3846} \\
Hit@10    & \num{0.3183} & \num{0.4085} & \num{0.4520} \\
Hit@20    & \num{0.3881} & \num{0.4789} & \num{0.5416} \\
Hit@50    & \num{0.4929} & \num{0.5920} & \num{0.6658} \\
Recall@1  & \num{0.1905} & \num{0.2026} & \num{0.2180} \\
Recall@5  & \num{0.2569} & \num{0.2999} & \num{0.3394} \\
Recall@10 & \num{0.3144} & \num{0.3651} & \num{0.4049} \\
Recall@20 & \num{0.3833} & \num{0.4316} & \num{0.4938} \\
Recall@50 & \num{0.4870} & \num{0.5478} & \num{0.6210} \\
Precision@10 & \num{0.0362} & \num{0.0482} & \num{0.0557} \\
MRR@10    & \num{0.2268} & \num{0.2797} & \num{0.3123} \\
MRR@50    & \num{0.2349} & \num{0.2882} & \num{0.3225} \\
nDCG@10   & \num{0.2476} & \num{0.2883} & \num{0.3214} \\
nDCG@50   & \num{0.2880} & \num{0.3319} & \num{0.3733} \\
\bottomrule
\end{tabular}
\end{table}

The ordering is consistent across every metric: My\_RAG first,
Metadata-RAG second, LongRAG third. This uniformity is itself a
finding: the three architectures rank differently along a single
dominant axis, and the axis is not noise.

\subsection{Relative position}
\begin{center}
\small
\begin{tabular}{l r}
\toprule
\textbf{Comparison} & \textbf{Relative delta} \\
\midrule
Metadata-RAG vs LongRAG & $+16.1\%$ on macro Recall@10 \\
Metadata-RAG vs My\_RAG & $-9.8\%$ on macro Recall@10 \\
My\_RAG vs LongRAG & $+28.8\%$ on macro Recall@10 \\
\bottomrule
\end{tabular}
\end{center}

\subsection{Per-dataset relative deltas}
\begin{table}[h]
\centering
\caption{Relative improvement of Metadata-RAG over LongRAG per
dataset.}
\label{tab:relativedelta}
\small
\begin{tabular}{l r r r}
\toprule
\textbf{Dataset} & \textbf{LongRAG} & \textbf{Metadata-RAG}
  & \textbf{Relative delta} \\
\midrule
CLERC              & \num{0.0786} & \num{0.2726} & $+247\%$ \\
CUAD               & \num{0.0324} & \num{0.0294} & $-9.3\%$ \\
GreekBarBench      & \num{0.2606} & \num{0.2711} & $+4.0\%$ \\
LaRA               & \num{0.8861} & \num{0.8874} & $+0.15\%$ \\
\bottomrule
\end{tabular}
\end{table}

The largest advantage of Metadata-RAG over LongRAG is on the
shortest-document corpus (CLERC), where LongRAG's coarse
multi-document units dilute the relevant document. The single
dataset on which Metadata-RAG loses to LongRAG is CUAD, where the
query text is templated and produces at most one constraint per
query, leaving the metadata filter inactive and the dense-only
Metadata-RAG pipeline at a slight disadvantage relative to LongRAG's
coarser but better-scored long units on long legal contracts.

% =====================================================================
\section{Analysis}\label{sec:analysis}

\subsection{Why My\_RAG wins}
My\_RAG combines three architectural mechanisms that neither of the
other two tracks provides simultaneously:

\begin{enumerate}[leftmargin=1.5em,itemsep=0.15em]
  \item hierarchical chunking that never crosses semantic-parent
    boundaries, which keeps evidence coherent at the retrieval unit
    level;
  \item a learned-sparse lexical signal fused with the dense signal
    by Reciprocal Rank Fusion, which recovers evidence that the
    dense signal alone misses;
  \item a parent-aggregation rule that rewards parents with multiple
    strong child chunks, which prevents a single long parent from
    dominating the document ranking by sheer chunk count.
\end{enumerate}

None of the other tracks has all three at once. LongRAG has none of
them; Metadata-RAG has the first, but only through amendment 04, and
its metadata filter is inactive on three of the four datasets, so
its effective architecture on those three datasets is a plain
dense-only document-level retriever.

\subsection{Why Metadata-RAG beats LongRAG}
The comparison is not between two dense-only retrievers at the same
granularity. LongRAG retrieves over units that may contain many
documents; Metadata-RAG retrieves over windows that are always
single-document. On corpora whose documents are shorter than
LongRAG's unit budget, the single-document window preserves
document identity in the ranking, while the multi-document unit
dilutes it. The magnitude of the resulting gap is the largest on
CLERC, where LongRAG packs roughly eleven documents per unit; it
narrows on LaRA, where both architectures operate at
near-document-level granularity.

\subsection{Why Metadata-RAG loses to My\_RAG}
On the three datasets where the metadata filter is inactive,
Metadata-RAG is a dense-only single-document retriever with no
hierarchical chunking, no lexical signal, and no parent aggregation.
My\_RAG has all three. The gap of \num{9.8}\% on the primary metric
is the aggregate effect of those missing mechanisms.

On LaRA, where the metadata filter does activate, the gap narrows to
\num{5.3}\%. This is consistent with the metadata filter contributing
a real but bounded improvement on the one corpus where it activates.

\subsection{The metadata-stage activation finding}
The single most important observation for the program's paper is the
metadata-stage activation rate. Under the frozen
\texttt{min\_match\_ratio} of \num{0.5}, the filter activated on
\num{0.1}\% of CLERC queries, \num{2.4}\% of CUAD, \num{7.7}\% of
GreekBarBench, and \num{66.3}\% of LaRA. The pattern is not random:
it tracks the constraint density produced by the frozen extractor
on each dataset's query texts. Long judicial opinions in CLERC emit
dozens of constraints per query and make the frozen threshold
unattainable. Templated CUAD queries emit at most one constraint.
Greek queries emit nothing because the extractor is ASCII-only.
English natural questions in LaRA emit a moderate number of
constraints and the filter activates on the majority of the corpus.

The architectural observation is that a metadata-driven retrieval
architecture with a frozen threshold can be effectively dense-only
on corpora whose query text distribution does not align with the
threshold. No post-hoc tuning was applied; the observation is
reported as-is.

% =====================================================================
\section{Threats to Validity}\label{sec:threats}

\subsection{Scope}
The comparison covers four of the program's canonical datasets. The
remaining three (RULER, Loong, LongBench-v2) are excluded by
explicit scoping decision to match the runnable scope of all three
tracks.

\subsection{Statistical uncertainty}
Both tracks are executed once, deterministically, on the full query
set. No bootstrap confidence intervals, no repeated trials, and no
paired significance tests are reported. The observed differences are
descriptive.

\subsection{Precision}
FP16 embeddings are used. FP32 equivalence is not claimed.

\subsection{Ground-truth projection}
All three tracks are evaluated at the canonical-document level.
Datasets with span-level or clause-level annotations could arguably
be evaluated at finer granularity; the program's shared protocol
commits to document level.

\subsection{Controlled adaptations}
Each track contains controlled adaptations. Their full documentation
is in each track's own report; the comparison reports only the
executed numbers.

\subsection{Frozen thresholds}
The metadata filter threshold of \num{0.5}, the fusion coefficients
$\alpha_m=0.3$ and $\alpha_s=0.7$, the RRF constant $k=60$, the
child-chunk size of \num{256} tokens, the parent-aggregation weight
$w=0.05$, and the long-unit budget of $\approx\num{4096}$ tokens are
all frozen. Any of them could in principle be tuned to improve a
given track's numbers on a given corpus; the controlled comparison
does not permit this.

% =====================================================================
\section{Conclusion}

We have compared LongRAG, Metadata-RAG, and My\_RAG under a fully
controlled, retrieval-only protocol on four heterogeneous
long-context corpora. Under this protocol My\_RAG achieves a macro
document \texttt{Recall@10} of \num{0.4049}, Metadata-RAG of
\num{0.3651}, and LongRAG of \num{0.3144}. The ordering is
consistent across every reported metric and every dataset. The
magnitude of the gaps decomposes cleanly along the architectural
axes: unit granularity (LongRAG's coarse multi-document units versus
Metadata-RAG's single-document windows versus My\_RAG's parent-bound
child chunks), retrieval signal (dense-only versus dense plus
metadata versus dense plus learned-sparse plus RRF), and aggregation
(first-occurrence versus parent-based best-plus-margin).

The most useful finding for the program's paper is the interaction
between the frozen metadata-filter threshold and the constraint
density of the corpora. Metadata-RAG is dense-only on three of the
four datasets and metadata-filtered only on LaRA. This is a
reproducible property of the frozen design and is reported without
post-hoc adjustment. It is the kind of architectural observation
that a controlled comparison is meant to produce.

\end{document}